\begin{exercise}{9}
Show that the datum of an adjunction in the sense of Definition~1.1.17
is equivalent to the datum of a pair of functors $(F,G)$ together with
natural transformations $\varepsilon\colon FG\to\mathrm{id}$ and
$\eta\colon\mathrm{id}\to GF$\footnote{The original prints
$\eta\colon GF\to\mathrm{id}$. The correct direction is
$\eta\colon\mathrm{id}\to GF$: its component at $X$ must be
$X\to GF(X)$ to induce the first arrow of the first displayed
composite below.} satisfying the triangle identities, i.e., the
obvious composites
\[
 F(X)\to F(GF(X))\cong FG(FX)\to F(X)
\]
and
\[
 G(X)\to GF(G(X))\cong G(FG(X))\to G(X)
\]
are the identity of $F(X)$ and $G(X)$, respectively.

\begin{proof}
Write $F\colon\mathcal C\to\mathcal D$ and
$G\colon\mathcal D\to\mathcal C$. Given the natural bijection of
Definition~1.1.17,
\[
 \Phi_{x,y}\colon\operatorname{Hom}_{\mathcal D}(Fx,y)
       \xrightarrow{\cong}\operatorname{Hom}_{\mathcal C}(x,Gy),
\]
set
$\eta_x=\Phi_{x,Fx}(\mathrm{id}_{Fx})$ and
$\varepsilon_y=\Phi^{-1}_{Gy,y}(\mathrm{id}_{Gy})$.
Naturality of $\Phi$ and its inverse gives
\begin{equation}\label{eq:adjunction-formulas}
 \Phi(f)=G(f)\eta_x,\qquad
 \Phi^{-1}(g)=\varepsilon_yF(g).
\end{equation}
For $a\colon x\to x'$ and $b\colon y\to y'$, naturality also gives
\[
 GF(a)\eta_x=\Phi_{x,Fx'}(F(a))=\eta_{x'}a,
 \qquad
 b\varepsilon_y=\Phi^{-1}_{Gy,y'}(G(b))
                  =\varepsilon_{y'}FG(b).
\]
Thus $\eta$ and $\varepsilon$ are natural transformations.
Applying \eqref{eq:adjunction-formulas} to
$\Phi^{-1}(\eta_x)=\mathrm{id}_{Fx}$ and
$\Phi(\varepsilon_y)=\mathrm{id}_{Gy}$ yields
\[
 \varepsilon_{Fx}F(\eta_x)=\mathrm{id}_{Fx},\qquad
 G(\varepsilon_y)\eta_{Gy}=\mathrm{id}_{Gy},
\]
the required triangle identities.

Conversely, suppose natural transformations $\eta$ and $\varepsilon$
with these identities are given. Define
$\Phi(f)=G(f)\eta_x$ and $\Psi(g)=\varepsilon_yF(g)$.
They are natural in $x$ and $y$, and
\begin{align*}
 \Psi\Phi(f)
   &=\varepsilon_yFG(f)F(\eta_x)
     =f\varepsilon_{Fx}F(\eta_x)=f,\\
 \Phi\Psi(g)
   &=G(\varepsilon_y)GF(g)\eta_x
     =G(\varepsilon_y)\eta_{Gy}g=g.
\end{align*}
The middle equalities use naturality of $\varepsilon$ and $\eta$,
respectively. Hence $\Phi$ is a natural bijection and defines an
adjunction. Formula~\eqref{eq:adjunction-formulas} recovers the original
adjunction bijection, while evaluating the constructed maps at
identities recovers the original $\eta$ and $\varepsilon$.
The two constructions are therefore inverse.
\end{proof}
\end{exercise}
